\section*{Bab \ref{ch:b3:extension-strike-slip} --- Tektonik Perpanjangan dan Geser-Mendatar}

\begin{solution}{exo:b3:extension-strike-slip:1}
Lapisan memiliki kemiringan $50 - 35 = \ang{15}$, menjauhi sesar; $\beta =
\sin\ang{50}/\sin\ang{35} = 0.766/0.574 = 1.34$.
\end{solution}

\begin{solution}{exo:b3:extension-strike-slip:2}
Langkah dekstral ke kanan: cekungan. Langkah dekstral ke kiri: punggungan. Langkah
sinistral ke kiri: cekungan. Langkah sinistral ke kanan: punggungan.
\end{solution}

\begin{solution}{exo:b3:extension-strike-slip:3}
$105/20 = \qty{5.25}{km/Ma} = \qty{5.25}{mm/yr}$, mendekati
\qty{6}{mm/yr}: laju itu tidak harus konstan selama \qty{20}{Ma}.
\end{solution}

\begin{solution}{exo:b3:extension-strike-slip:4}
\qty{200}{\percent} berarti $\Delta L = 2L_0$: $L_0 = \qty{117.5}{km}$ dan $L
= \qty{352.5}{km}$. Di utara, $\Delta L = 0.5L_0$: $L_0 = \qty{470}{km}$
dan $L = \qty{705}{km}$.
\end{solution}

\begin{solution}{exo:b3:extension-strike-slip:5}
$563/20 = 28$ sampai $563/15 = \qty{38}{mm/yr}$. Pergeseran itu merupakan nilai
minimum, dan pergerakan yang diukur dengan survei melintasi zona sesar juga
mencakup gerak pada sesar-sesar tetangga dalam sistem dan regangan elastis yang
tersimpan di antara gempa bumi.
\end{solution}

\begin{solution}{exo:b3:extension-strike-slip:6}
$t = 90 - \text{kemiringan}$. Pada \ang{60}: $(1 + 0.6\times1.732)/(1 - 0.6\times
0.577) = 2.039/0.654 = 3.12$; pada \ang{50}: $1.715/0.497 = 3.45$; pada
\ang{40}: $1.504/0.285 = 5.28$. Penguncian terjadi pada $\arctan0.6 = \ang{31.0}$.
\end{solution}

\begin{solution}{exo:b3:extension-strike-slip:7}
$\arctan0.4 = \ang{21.8}$ dan $\arctan0.2 = \ang{11.3}$. Sesar yang miring
\ang{15} hanya tergelincir jika $\mu < \tan\ang{15} = 0.27$.
\end{solution}

\begin{solution}{exo:b3:extension-strike-slip:8}
$2.5\times4.0 = \qty{10.0}{km^2}$ dan $2.5\times(4.0 - 1.5) =
\qty{6.25}{km^2}$.
\end{solution}

\begin{solution}{exo:b3:extension-strike-slip:9}
$132/1.19 = \qty{110.9}{km}$, yang meregang sebesar \qty{21.1}{km}, dan itu
membutuhkan $21.1/6 = \qty{3.5}{Ma}$.
\end{solution}

\begin{solution}{exo:b3:extension-strike-slip:10}
$3.0/\sin\ang{60} = \qty{3.464}{km}$ sebelum dan $3.0/\sin\ang{45} =
\qty{4.243}{km}$ sesudah: \qty{0.779}{km} per blok ($\beta = 1.22$). Lapisan
memiliki kemiringan $60 - 45 = \ang{15}$. Setiap blok berputar secara utuh:
rotasi yang menurunkan kemiringan sesar yang, misalnya, miring ke timur
memutar garis horizontal di dalam blok sehingga naik ke timur, dan lapisan
miring ke barat, berlawanan dengan sesarnya.
\end{solution}

\begin{solution}{exo:b3:extension-strike-slip:11}
Pada zona timur--barat sinistral $\sigma_1$ berarah \ang{045}. Sesar R berada
pada $\phi/2 = \ang{20}$ dari zona di sisi $\sigma_1$: jurus \ang{070},
sinistral. R$'$ pada $90 - 20 = \ang{70}$: jurus \ang{020}, dekstral. Rekahan
tarik sejajar dengan $\sigma_1$: jurus \ang{045}.
\end{solution}

\begin{solution}{exo:b3:extension-strike-slip:12}
Untuk $\mu = 0.6$ sesar terbentuk pada \ang{60.5} dan terkunci pada \ang{31.0}: $\beta =
\sin\ang{60.5}/\sin\ang{31.0} = 0.870/0.515 = 1.69$. Untuk $\mu = 0.85$: dari
\ang{65.2} sampai \ang{40.4}, $\beta = 0.908/0.648 = 1.40$. Sesar dengan gesekan
lebih rendah berotasi lebih jauh sebelum terkunci: batuan lemah meregang lebih
banyak sebelum sesar baru harus terbentuk.
\end{solution}

\begin{solution}{pb:b3:extension-strike-slip:1}
\textbf{1.} Lapisan semula horizontal dan telah berputar bersama blok; sesar
juga berputar dengan sudut yang sama, sehingga $\theta_0 = 42 + 18 = \ang{60}$,
yaitu kemiringan Anderson untuk $\phi = \ang{30}$.

\textbf{2.} $\beta = \sin\ang{60}/\sin\ang{42} = 0.8660/0.6691 = 1.294$.

\textbf{3.} $24.0/1.294 = \qty{18.5}{km}$; domain telah meregang sebesar
\qty{5.5}{km} (\qty{5.46}{km}).

\textbf{4.} $t = \ang{48}$: $(1 + 0.6\times0.900)/(1 - 0.6\times1.111) =
1.540/0.334 = 4.62$, dibandingkan $3.12$ untuk sesar baru pada \ang{60.5}: kira-kira
separuh lebih banyak.

\textbf{5.} Pada $\arctan0.6 = \ang{31.0}$; $\beta = \sin\ang{60}/\sin\ang{31}
= 1.68$ paling besar. Sebelum itu, sesar baru pada sekitar \ang{60} akan
memotong blok, dan generasi kedua dari domino akan mengambil alih.

\textbf{6.} Setiap kolom dinding gantung turun sebesar penurunan sesar di
bawahnya sepanjang lompatan horizontal; sesar paling curam di dekat permukaan
dan melandai menjauhinya, sehingga kolom yang dekat dengan sesar turun paling
banyak dan yang jauh turun paling sedikit: lapisan miring ke arah sesar.

\textbf{7.} $D = A/H = 16.0/3.2 = \qty{5.0}{km}$.

\textbf{8.} $D$ sebanding dengan $A$: $\qty{5.0}{km} \pm \qty{0.5}{km}$.

\textbf{9.} $3.2\times(5.0 - 2.0) = \qty{9.6}{km^2}$.

\textbf{10.} Lompatan horizontal, \qty{3.2}{km}.

\textbf{11.} Isian menebal ke arah sesar, membentuk baji, dan lapisan yang lebih
tua lebih miring daripada lapisan yang lebih muda: setiap lapisan diendapkan
ketika rollover dan ruang di samping sesar terus bertambah.

\textbf{12.} $z = P/(\rho g) = \num{5.5e8}/(2800\times9.8) = \qty{20.0}{km}$.

\textbf{13.} $20.0/(24 - 16) = \qty{2.5}{km/Ma} = \qty{2.5}{mm/yr}$.

\textbf{14.} $20.0/\sin\ang{30} = \qty{40}{km}$ dari pergeseran, dan $40\cos\ang{30} =
\qty{34.6}{km}$ dari perpanjangan horizontal.

\textbf{15.} Tidak: \ang{15} berada di bawah kemiringan kunci \ang{31}. Entah sesar
itu bergeser ketika lebih curam lalu terputar (\omterm{def:b3:extension-strike-slip:core-complex}{engsel bergulir}), atau sesar
itu lemah, dengan $\mu < 0.27$ akibat lempung atau fluida bertekanan tinggi.

\textbf{16.} Dengan $45 - 15 = \ang{30}$: ketika lempeng atas bergerak menjauh,
dinding bawah yang tidak lagi terbebani naik karena isostasi dan melengkung,
sehingga sesar terputar melalui engsel.

\textbf{17.} \qty{34.6}{km} dibandingkan \qty{8.7}{km} milik paruh-graben (pertanyaan
23): empat kali lebih besar. \omterm{def:b3:extension-strike-slip:core-complex}{Kompleks inti metamorf} terletak di tempat
perpanjangan paling besar, seperti di Basin and Range bagian tengah.

\textbf{18.} Sesar sinistral yang melangkah ke kiri: pelepasan, cekungan terbuka.

\textbf{19.} $\qty{5}{mm/yr}\times\qty{1.2}{Ma} = \qty{6.0}{km}$.

\textbf{20.} $4.0\times6.0 = \qty{24}{km^2}$; $3.0 + 6.0 = \qty{9.0}{km}$ panjangnya;
$9.0/4.0 = 2.25$.

\textbf{21.} $132/10 = 13$ sampai $132/7 = 19$: cekungan Laut Mati jauh lebih
memanjang, pergeserannya beberapa kali lipat lebarnya.

\textbf{22.} \omterm{def:b3:extension-strike-slip:bends}{Tekukan penahan}: punggungan terdorong ke atas di antara untaian; dalam
penampang, \omterm{def:b3:extension-strike-slip:flower}{struktur bunga} positif, sesar bercabang ke atas menjadi untaian
dengan komponen naik, seperti yang selebar 100 sampai \qty{300}{m} pada Transform
Laut Mati.

\textbf{23.} $3.2 + 5.46 = \qty{8.66}{km}$: paruh-graben Merrow telah
diperpanjang sebesar \textbf{\qty[text-series-to-math=true]{8.7}{km}}.
\end{solution}
